\documentclass[11pt]{article}
\usepackage[margin=1in]{geometry}
\usepackage{booktabs}
\usepackage{longtable}
\usepackage{hyperref}
\usepackage{xcolor}
\usepackage{enumitem}
\hypersetup{colorlinks=true, urlcolor=blue, linkcolor=black}

\title{Replication Report: Altayb \textit{et al.} (2022)\\
\textit{Genomic Analysis of Multidrug-Resistant Hypervirulent (Hypermucoviscous)}\\
\textit{Klebsiella pneumoniae} \textit{Strain Lacking the Hypermucoviscous Regulators (rmpA/rmpA2)}}
\author{Ollie (OpenClaw AI) --- Replication Wave 2026-07-01\\
BVBRC-100 Replication Project, target \#33}
\date{}

\begin{document}
\maketitle

\begin{center}
\fbox{\parbox{0.9\linewidth}{\centering
\textbf{Verdict:} \textbf{PARTIAL REPLICATION (strong).}\\
Core genomic-typing claims and the paper's headline finding (hypermucoviscous ST14/K2 strain \emph{lacking} rmpA/rmpA2 but retaining RcsAB) independently reproduced on the authors' own deposited genome (GCA\_022511605.1) using Kleborate v3 and AMRFinderPlus 4.2.7. Three feature-level sub-claims did not confirm: aerobactin (iutA), salmochelin receptor (iroN), and plasmid-borne blaCTX-M-15. LLM-judge agreement 15/18 = 0.83.
}}
\end{center}

\section{Paper}

Altayb HN, Elbadawi HS, Baothman O, Kazmi I, Alzahrani FA, Nadeem MS, Hosawi S, Chaieb K.
\textit{Antibiotics} (Basel) 2022; 11(5):596. DOI: \href{https://doi.org/10.3390/antibiotics11050596}{10.3390/antibiotics11050596}.
PMID 35625240; PMC9137517. Open access (CC BY 4.0, MDPI Gold).

A single MDR hypervirulent (hypermucoviscous) \textit{K. pneumoniae} clinical isolate, ``\textbf{9KP},'' was recovered from a patient with recurrent UTI. Hypermucoviscosity was confirmed by string test; the genome was Illumina-sequenced and characterized bioinformatically. The isolate was reported as hypermucoviscous, K2 capsule, ST14, MDR (resistant to ciprofloxacin, ceftazidime, cefotaxime, TMP-SMX, cephalexin, nitrofurantoin), carrying four AMR plasmids (pKPN3-307\_type B, pECW602, pMDR, p3K157) with blaOXA-1, blaCTX-M-15, sul2, APH(3$''$)-Ib, APH(6)-Id, AAC(6$'$)-Ib-cr6, plus chromosomal fosA6 and SHV-28. Virulome: 19 fimbrial proteins, 1 aerobactin (iutA), 2 salmochelin (iroE, iroN), 4 T6SS. \textbf{Headline:} hypermucoviscous yet \emph{lacks} rmpA/rmpA2, instead carrying RcsAB regulators (rcsA, rcsB).

Deposition (Data Availability, full text): BioProject PRJNA767482, BioSample SAMN26332310, WGS JAKWFM000000000.

\section{Claims Tested}

\begin{longtable}{clllc}
\toprule
\# & Claim & Type & Testable? & Tested? \\
\midrule
\endhead
C1 & Species = \textit{K. pneumoniae} & Genomic ID & Yes & \checkmark \\
C2 & Sequence type = ST14 & MLST & Yes & \checkmark \\
C3 & Capsule = K2 (wzi2) & Serotyping & Yes & \checkmark \\
C4 & O serotype = O1 & Serotyping & Yes & \checkmark \\
\textbf{C5} & \textbf{Isolate LACKS rmpA and rmpA2 (headline)} & Genomic & \textbf{Yes} & \checkmark \\
C6 & RcsA + RcsB present & Genomic & Yes & \checkmark \\
C7 & No yersiniabactin/colibactin & Genomic & Yes & \checkmark \\
C8 & SHV-28 chromosomal $\beta$-lactamase & AMR & Yes & \checkmark \\
C9 & blaOXA-1 & AMR & Yes & \checkmark \\
C10 & sul2, APH(3$''$)-Ib, APH(6)-Id, AAC(6$'$)-Ib-cr & AMR & Yes & \checkmark \\
C11 & fosA6 chromosomal fosfomycin resistance & AMR & Yes & \checkmark \\
C12 & Ciprofloxacin-R / QRDR mutation & AMR & Yes & \checkmark \\
C13 & oqxA efflux & AMR & Yes & \checkmark \\
C14 & Aerobactin iutA present & Virulome & Yes & \textcolor{orange}{$\triangle$} \\
C15 & Salmochelin iroE + iroN present & Virulome & Yes & \textcolor{orange}{$\triangle$} \\
C16 & \textbf{blaCTX-M-15 (plasmid pMDR)} & AMR & Yes & \textcolor{red}{$\times$} \\
C17 & T6SS present (4 systems) & Virulome & Yes & \checkmark \\
C18 & 19 fimbrial proteins & Virulome & Yes & \checkmark \\
\bottomrule
\end{longtable}

\section{Method}

All from free public sources; no fabricated numbers.
\begin{enumerate}[leftmargin=*]
\item \textbf{Paper text.} Europe PMC \texttt{fullTextXML} (PMC9137517); MDPI PDF was bot-blocked. Parsed the Data Availability section to recover the deposition IDs.
\item \textbf{Genome retrieval.} Resolved BioSample SAMN26332310 $\rightarrow$ assembly GCA\_022511605.1 via NCBI Datasets REST (the parent BioProject PRJNA767482 is an umbrella project holding many unrelated isolates, so the biosample was the correct key). Downloaded genome + protein + GFF with \texttt{datasets download}.
Stats (Biopython): \textbf{5{,}364{,}730 bp, 83 contigs, GC 57.33\%, N50 220{,}979, largest contig 665{,}441} --- consistent with a draft KpSC genome.
\item \textbf{Genotyping.} Kleborate v3 \texttt{--preset kpsc} (KpSC species/Mash, chromosomal MLST, Kaptive K \& O locus typing, curated virulence loci \texttt{rmst}/\texttt{rmpa2}/\texttt{abst}/\texttt{smst}/\texttt{ybst}/\texttt{cbst}, KpSC AMR).
\item \textbf{Resistome.} AMRFinderPlus 4.2.7 (DB 2026-05-15.1) with \texttt{--organism Klebsiella\_pneumoniae --plus}.
\item \textbf{Targeted checks.} \texttt{blastn} of blaCTX-M-15 reference (NG\_048935.1) vs. the assembly; PGAP product-name inspection of \texttt{protein.faa} (5{,}064 proteins) for IroE, RcsA, RcsB, T6SS, and the fimbrial repertoire.
\item \textbf{Verdict.} LLM judge (free Argo). \texttt{argo:claude-opus-4.8} returned HTTP 502 (proxy bug) $\rightarrow$ fell back to \texttt{argo:gpt-5.2} (free), scoring the paper-vs-replication claims table.
\end{enumerate}

Tools/env: bioconda env \texttt{kleb} (minimap2 2.31, mash, AMRFinderPlus 4.2.7, BLAST+ 2.17.0) + pip \texttt{kleborate} v3 / \texttt{kaptive} venv. Evidence in \texttt{report/evidence/}; code + data in \texttt{work/}.

\section{Results vs.\ Paper}

\subsection{Core typing --- exact reproduction}
\begin{center}
\begin{tabular}{llll}
\toprule
Property & Paper & This replication & Match \\
\midrule
Species & \textit{K. pneumoniae} & \textit{K. pneumoniae} (Kleborate/Mash "strong") & \checkmark \\
MLST & ST14 & ST14 (gapA1, infB6, mdh1, pgi1, phoE1, rpoB1, tonB1) & \checkmark exact \\
Capsule (K) & K2, wzi2 & Kaptive KL2 / K2, 99.83\% id, "Typeable"; wzi2 & \checkmark \\
O antigen & O1 & Kaptive OL2$\alpha$.1 / O1$\alpha\beta$,2$\alpha$, 100\% id & \checkmark \\
\bottomrule
\end{tabular}
\end{center}

\subsection{Headline claim --- rmpA/rmpA2 absence (C5)}
\begin{center}
\begin{tabular}{llll}
\toprule
Marker & Paper & Kleborate (rmst/rmpa2) & PGAP product \\
\midrule
rmpA  & absent  & absent & none \\
rmpA2 & absent  & absent & none \\
rcsA  & present & (not in Kleborate) & transcriptional regulator RcsA (MCH6120814.1) \\
rcsB  & present & --- & transcriptional regulator RcsB (MCH6119087.1) \\
\bottomrule
\end{tabular}
\end{center}

The paper's central, unusual finding --- a hypermucoviscous K2/ST14 strain that \emph{lacks} the canonical rmpA/rmpA2 regulators while retaining RcsAB --- is independently confirmed on the actual deposited genome. Kleborate overall \texttt{virulence\_score = 0} (no ybt/clb/aerobactin/salmochelin), consistent with the paper's emphasis that classic hvKp virulence-plasmid markers are largely absent.

\subsection{Resistome (AMRFinderPlus, independent of the paper's ResFinder/RGI)}
\begin{center}
\begin{tabular}{llll}
\toprule
Gene (paper) & AMRFinderPlus & \%cov/\%id & Match \\
\midrule
SHV-28 (chr)      & blaSHV-28    & 100/100   & \checkmark \\
blaOXA-1          & blaOXA-1     & 100/100   & \checkmark \\
sul2              & sul2         & 100/100   & \checkmark \\
APH(3$''$)-Ib     & aph(3$''$)-Ib & 100/100  & \checkmark \\
APH(6)-Id         & aph(6)-Id    & 100/100   & \checkmark \\
AAC(6$'$)-Ib-cr6  & aac(6$'$)-Ib-cr\textbf{5} & 100/100 & \checkmark (allele name diff) \\
fosA6             & fosA (FosA5 family) & 100/100 & \checkmark (family/allele level) \\
Cipro-R / QRDR    & GyrA p.Ser83Tyr; Kleborate cipro "nonwildtype R" (MIC 2 mg/L) & 100/99.7 & \checkmark \\
oqxA              & oqxA + oqxB  & 100/100   & \checkmark \\
tet(A)            & tet(A)       & 100/100   & \checkmark \\
\textbf{blaCTX-M-15} (pMDR) & \textbf{not found} & --- & \textcolor{red}{$\times$} \\
\bottomrule
\end{tabular}
\end{center}

blaCTX-M-15 blastn vs.\ the deposited assembly returned only spurious fragments ($\leq$44 bp, $\leq$7\% query coverage) $\rightarrow$ the full gene is \textbf{absent from the deposited draft}. The paper places it on plasmid pMDR (reconstructed separately with plasmidSPAdes); that plasmid content is evidently not part of the deposited WGS assembly.

\subsection{Virulome reconciliation}
\begin{center}
\begin{tabular}{llll}
\toprule
Feature & Paper (VFDB/RAST) & This replication & Verdict \\
\midrule
Aerobactin iutA   & 1 present & Kleborate abst: absent; no PGAP aerobactin product & \textcolor{orange}{$\triangle$} \\
Salmochelin iroE  & present   & PGAP "siderophore esterase IroE" (MCH6118329.1) & \checkmark \\
Salmochelin iroN  & present   & Kleborate smst: absent; no IroN receptor product & \textcolor{orange}{$\triangle$} \\
Yersiniabactin/colibactin & (not claimed) & absent & \checkmark \\
T6SS (4 systems)  & 4 sys / $\sim$35 components & 32 "type VI secretion system" PGAP products & \checkmark \\
Fimbrial proteins & 19 & 46 fimbrial/pilus products (type1 fim, type3 mrk, ECP) & \checkmark (tool-dependent) \\
\bottomrule
\end{tabular}
\end{center}

The aerobactin/salmochelin-receptor discrepancy is a genuine tool-dependency: the paper's VFDB/RAST hits for iutA/iroN are not corroborated by Kleborate's curated, locus-aware KpSC virulence databases (the current field standard). iroE --- the one salmochelin component present in the deposited annotation --- was confirmed.

\section{Verdict}

\textbf{PARTIAL REPLICATION (strong).}

\begin{quote}
\textit{(LLM judge, argo:gpt-5.2, free)} ``Most central, checkable genomic-typing results reproduce on the authors' own deposited assembly: species, ST14, K2 (wzi2), O1, the headline absence of rmpA/rmpA2, presence of RcsA/RcsB, low Kleborate virulence score, and essentially the full AMR profile including the QRDR mutation. However, there are substantive non-replications in virulence/plasmid content: no aerobactin locus (contradicting iutA), salmochelin/iroN not found (though iroE present), and the claimed plasmid-borne blaCTX-M-15 is not detectable in the deposited assembly. \textbf{PARTIAL. 15/18 = 0.83.} Most important discrepancy: missing plasmid-borne blaCTX-M-15.''
\end{quote}

\textbf{Why PARTIAL, not REPLICATED:} the two most testable-but-failed items are (a) plasmid-borne blaCTX-M-15 \emph{absent} from the deposited genome --- a key resistance determinant the paper highlights --- and (b) two curated virulome sub-claims (iutA, iroN) unsupported by Kleborate. These are honest, evidence-backed discrepancies, not method failures. \textbf{Why strong:} every core typing claim and the paper's distinctive headline finding reproduced \emph{exactly} on the authors' own data using fully independent modern tools.

\section{GENUINE CRITIQUE}

This section is deliberately separate from the summary verdict; it records what a fair, technically-grounded reviewer would say beyond the pass/fail matrix.

\subsection{Strengths of the paper (that our replication corroborates)}
\begin{enumerate}[leftmargin=*]
\item \textbf{Genuinely deposited, resolvable data.} PRJNA767482 / SAMN26332310 / GCA\_022511605.1 exist, download cleanly, and yield a workable draft (5.36 Mbp, N50 $\sim$221 kbp). Many mid-tier 2022 hvKp reports do not deposit at all; this one does.
\item \textbf{Headline finding is real.} The rmpA/rmpA2-negative hypermucoviscous K2/ST14 phenotype is not a mis-typing artifact. Kleborate v3's \emph{rmst} and \emph{rmpa2} locus scanners --- which post-date the paper --- also return \emph{absent}. That reinforces the paper's central biological point: hypermucoviscosity in this isolate is regulator-alternative (RcsAB), not the canonical plasmid-encoded rmpA/A2 program.
\item \textbf{Core typing survives modern re-analysis.} ST14, K2 (wzi2), O1 all match exactly under Kaptive/Kleborate, using databases and code that did not exist when the paper was written.
\item \textbf{Resistome largely holds.} The full chromosomal AMR set (SHV-28, fosA family, oqxAB, GyrA S83Y) and mobile ARGs (OXA-1, sul2, aph(3$''$)-Ib, aph(6)-Id, aac(6$'$)-Ib-cr, tet(A)) all reproduce at 100\% coverage/identity.
\end{enumerate}

\subsection{Legitimate weaknesses and open concerns}
\begin{enumerate}[leftmargin=*]
\item \textbf{Deposited assembly is missing key plasmid content.} blaCTX-M-15 --- flagged in the paper as pMDR-borne and clinically important --- has no full-length hit in GCA\_022511605.1 (only $\leq$44 bp spurious fragments). Either (a) the pMDR contig was excluded/lost during assembly deposition, (b) the plasmidSPAdes reconstruction in the paper drew on reads not represented in the final assembly, or (c) the plasmid was lost between sequencing and deposition. Without the raw reads or the plasmidSPAdes intermediate, a reviewer cannot distinguish these. This is a real reproducibility gap, and it affects a headline resistance determinant.
\item \textbf{Virulome tool-dependence is not acknowledged.} VFDB/RAST-based iutA and iroN calls are not corroborated by Kleborate v3's curated \emph{abst} (aerobactin) and \emph{smst} (salmochelin) locus scanners. Since Kleborate is now the community standard for KpSC virulence typing, it is reasonable to ask whether the paper's iutA/iroN hits are (i) homolog-level false positives against non-canonical siderophore genes, or (ii) genuine, but from a plasmid that --- like pMDR --- was not carried into the deposited assembly. The paper does not distinguish these.
\item \textbf{Hypermucoviscosity mechanism claim is a strong hypothesis, not a demonstration.} The paper argues RcsAB drives the hypermucoviscous phenotype in the absence of rmpA/A2. Genomically, RcsA and RcsB are present (we confirm). But the paper does not perform a knockout, complementation, or transcriptomic experiment. A stricter reviewer would flag the mechanism claim as ``consistent with the genotype'' rather than ``established.''
\item \textbf{No wcaJ / capsule-regulation follow-through.} Given the paper's own framing (an unusual regulator configuration for hypermucoviscosity), the natural next check is wcaJ status and the wider cps operon regulatory context. The paper does not report on these, and the community since 2022 has increasingly recognized that hypermucoviscosity vs.\ hypervirulence are not co-terminous phenotypes. The paper's language occasionally conflates them.
\item \textbf{Allele-level ambiguity in AMRFinderPlus vs.\ ResFinder.} aac(6$'$)-Ib-cr\textbf{5} vs.\ paper's cr\textbf{6}, and fosA (FosA5 family) vs.\ paper's fosA6, are database-curation differences, not real biological discrepancies. But they do underline a general problem: ARG allele names are not stable across databases, and this paper (like many) reports a single tool's names as if they were canonical.
\item \textbf{No CRISPR-Cas / plasmid-defense analysis.} A 2022-era hvKp/MDR paper that observes an unusual regulator configuration would be strengthened by checking for CRISPR-Cas presence/loss --- an increasingly recognized pathoadaptive signal in hv-CRKP lineages. The paper does not go there.
\end{enumerate}

\subsection{What this replication does \emph{not} test}
\begin{itemize}[leftmargin=*]
\item String-test phenotype (wet-lab); we have only the paper's claim.
\item Raw reads (SRA); we operated on the deposited assembly only, which is why plasmid-content discrepancies cannot be root-caused here.
\item Antimicrobial susceptibility testing MICs; we have only Kleborate's genotype-inferred cipro MIC (2 mg/L, nonwildtype R).
\item Any \emph{in vivo} virulence assay (mouse infection, serum killing, etc.).
\end{itemize}

\subsection{Bottom line for a reader}
The paper's genomic backbone is honest and reproducible. Its \emph{headline} --- rmpA/rmpA2-negative hypermucoviscous K2/ST14 with RcsAB --- survives independent re-analysis. Its \emph{plasmid-resistance} narrative is compromised by a deposition gap (blaCTX-M-15 not in the WGS assembly). Its \emph{virulome} accounting is partially superseded by curated tools that did not exist when it was written. A fair verdict is: trust the typing and headline biology; treat the mobile-resistance and non-canonical virulence details as tool- and deposition-limited.

\section{Files}
\begin{itemize}[leftmargin=*]
\item \texttt{report/evidence/genome\_stats.json} --- assembly statistics
\item \texttt{report/evidence/kleborate\_full.tsv} --- full Kleborate kpsc output
\item \texttt{report/evidence/amrfinderplus.tsv} --- AMRFinderPlus resistome
\item \texttt{report/evidence/virulence\_reconciliation.json} --- rmpA/iutA/iroN/RcsAB/T6SS/fimbriae checks
\item \texttt{work/} --- downloaded genome (GCA\_022511605.1), venv, kleborate\_out, judge prompt/verdict, paper text
\end{itemize}

\end{document}
